% ==================================================================
% id:       lim-021
% topic:    גבולות ורציפות של פונקציות
% subtopic: רציפות ונקודות אי־רציפות
% solution: full
% ==================================================================
תהי
$f$
פונקציה
המוגדרת בסביבה של
$x_0$.
\begin{itemize}
    \item
    $x_0$
    נקראת נקודת אי-רציפות
    \underline{סליקה}
    של
    $f$,
    אם
    קיים
    $\displaystyle\lim_{x\to x_0}f(x)$,
    אך
    $\displaystyle\lim_{x\to x_0}f(x)\neq f(x_0)$.
    \item
    $x_0$
    נקראת אי-רציפות
    מסוג \underline{קפיצה}
    של
    $f$,
    אם
    הגבולות החד-צדדיים קיימים אך
    שונים, כלומר
    $\displaystyle\lim_{x\to x_0^-}f(x)\neq\displaystyle\lim_{x\to x_0^+}f(x)$.
    \item
    בכל מקרה אחר נגיד כי
    $x_0$
    נקראת אי-רציפות
    מסוג
    \underline{עיקרית} של
    $f$.
\end{itemize}

\textbf{
בכל אחד מהסעיפים הבאים,
מצאו וסווגו את כל נקודות האי-רציפות של הפונקציה (סליקה, קפיצה או עיקרית).
}
\bE
\item $\ds{f(x)=
\case{ccc}{\frac{x^2-9}{x-3} & x\neq 3 \\ 5 & x=3}}$
\begin{solbox}
בסביבת כל
$x_0\neq 3$
הפונקציה
$\frac{x^2-9}{x-3}$
היא אלמנטרית ולכן רציפה ב-
$x_0$.
נציג
\begin{align*}
    \frac{x^2-9}{x-3}&=\ \frac{(x-3)(x+3)}{x-3}\\
    &=\ x+3,
\end{align*}
מכאן
$$,\lim_{x\to3}f(x)=\lim_{x\to3}(x+3)=6\neq5=f(3)$$
נסיק כי
$f$
אינה רציפה ב-
$x=3$.
זאת נקודת אי-רציפות סליקה.
הערה: בחישוב הגבול השתמשנו באריתמטיקה של גבולות: הגבול של סכום, של מכפלה ושל מנה של פונקציות מתכנסות שווה לסכום, למכפלה ולמנה של הגבולות, כל עוד הגבול של המכנה שונה מ-$0$.
\end{solbox}
\item $\ds{f(x)=\case{ccc}{x & \abs{x}\leq1 \\ 1 & \abs{x}>1}}$
\begin{solbox}
נסמן
$$.g(x)=x,\ h(x)=1$$
שתי הפונקציות האלה רציפות בכל הציר הממשי. לכן עלינו לבדוק האם
$f$
רציפה בנקודות תפר, דהיינו ב-
$x=1$
ו-
$x=-1$.

נבדוק את
$x=1$.
מתקיים
\begin{align*}
    \lim_{x\to1^-}f(x)= \lim_{x\to1^-}x
    = 1
    =\lim_{x\to1^+}f(x),
\end{align*}
וידאנו עתה כי קיים ל-
$f$
גבול בנקודה
$x=1$
והוא שווה ל-
$1$,
הערך של הפונקציה בנקודה
$x=1$.
נסיק כי
$f$
רציפה ב-
$x=1$.

נבדוק את
$x=-1$.
מתקיים
\begin{align*}
    \lim_{x\to-1^-}f(x)= 1
    \neq\ -1
    = \lim_{x\to-1^+}x
    =\lim_{x\to-1^+}f(x),
\end{align*}
רואים כי  לא קיים ל-
$f$
גבול בנקודה
$x=-1$ -
הגבולות החד-צדדיים קיימים אך שונים.
מכאן נסיק כי
$f$
לא רציפה ב-
$x=-1$,
סוג האי-רציפות הוא קפיצה.
\end{solbox}
\item $\ds{f(x)=\case{ccc}{x\cdot\sin\lr{\frac1x} & x\neq 0 \\ 0 & x=0}}$
\hint{כדי לחשב את הגבול ב-$0$ היעזרו בכך ש-$\sin$ חסומה.}
\begin{solbox}
הפונקציה
$g(x)=x\cdot\sin(\tfrac{1}{x})$
רציפה לכל
$x\neq0$
כי היא אלמנטרית.
נשים לב כי
$$\lim_{x\to0}x\cdot\sin(\tfrac{1}{x})=0$$
כמכפלה של פונקציה אפסה בפונקציה חסומה.
מכאן רואים כי
$$,\lim_{x\to0}f(x)=0=f(0)$$
משמעו,
$f$
רציפה גם ב-
$x=0$,
אז ל-
$f$
אין נקודות אי-רציפות.
\end{solbox}
\item $\ds{f(x)=\case{ccc}{\frac1{1+e^{\frac1{x-1}}} & x\neq 1 \\ 1 & x=1}}$
\hint{חשבו בנפרד את הגבולות החד-צדדיים ב-$1$: לאן שואף $\frac1{x-1}$ מכל צד?}
\begin{solbox}
נשים לב כי
\begin{align*}
    \lim_{x\to1^-}f(x)= \lim_{x\to1^-}\frac1{1+e^{\frac1{x-1}}}
    \overset{(\ast)}{=} 1,
\end{align*}
$(\ast)$
שימו לב כי
$\displaystyle\lim_{t\to-\infty}e^t=0$.

בנוסף
\begin{align*}
    \lim_{x\to1^+}f(x)= \lim_{x\to1^+}\frac1{1+e^{\frac1{x-1}}}
    \overset{(\ast)}{=} 0.
\end{align*}
$(\ast)$
שימו לב כי
$\displaystyle\lim_{t\to\infty}e^t=\infty$.

אנו רואים כי הגבולות החד-צדדיים ב-
$x=1$
קיימים אך שונים,
לכן
$x=1$
היא נקודת אי-רציפות מסוג קפיצה.
אין לפונקציה הנתונה נקודות אי-רציפות אחרות,
כי בסביבת כל
$x_0\neq 1$
היא אלמנטרית ורציפה.
\end{solbox}
\eE
